\section{Constructing the local polynomial sections}
\label{sec:local-sections}
We now construct the fixed section and its price used in
Section~\ref{sec:local-splitting}. The construction uses a finite
collection of based, transported plaquettes. The finite collection
and all constants are chosen before the lattice period and threshold.

\subsection{A finite sampler for radial jets}
Let $\mathcal J_d$ be the finite-dimensional real space of radial
connection jets $A=(A_0,\ldots,A_{d-2})$. Here $A_j$ is a homogeneous
one-form of polynomial degree $j+1$ satisfying the radial identity
$\sum_i x_iA_{j,i}(x)=0$. It is this radial description of curvature
jets that underlies the label space. If $d<2$, the space is zero.
For a lattice path from the origin to $x$ and an oriented plaquette
at $x$, let $s(U)$ be the Lie-algebra curvature of that plaquette,
transported back to the origin. Its value is zero at the identity.
Define the sampler $S_s(A)$ by evaluating $s$ on the test holonomy
of the truncated radial connection $A$ at scale one.

The endpoint argument in Section~\ref{sec:local-splitting} shows
that $S_s$ is analytic near zero. Differentiating the link integral
equation at the zero connection gives the line integral of $A$.
At the identity, the derivative of a product is the sum of its
factor derivatives, and inversion changes their signs. The derivative
of transported plaquette curvature is consequently
\[
 DS_s(0)A=\oint_{\partial\square_{x,ij}}
                          \sum_{l=0}^{d-2}A_l.
\]
Transport along the based path contributes nothing to this derivative:
it conjugates a curvature whose value at zero is zero. The Lie
projection of the linearized plaquette is its Lie-algebra value.

These linear samplers jointly separate $\mathcal J_d$. To see this,
for a homogeneous $A_l$ write
\[
 N_l(x,h)=\int_0^1\!\int_0^1
   \partial_i A_{l,j}(x+hs e_j+hr e_i)\,dr\,ds
 -\int_0^1\!\int_0^1
   \partial_j A_{l,i}(x+hs e_i+hr e_j)\,dr\,ds.
\]
The fundamental theorem of calculus identifies the unit-square
circulation with $N_l(x,1)$. Also
$N_l(hx,h)=h^lN_l(x,1)$ and
$N_l(x,0)=\partial_i A_{l,j}(x)-\partial_j A_{l,i}(x)$.
Suppose all sampled circulations vanish. If $k$ is the largest
remaining degree, multiply the identity at an integer point $z$
by $h^k$ and use homogeneity to get
\[
                       \sum_{l=0}^k h^{k-l}N_l(hz,h)=0.
\]
For any real $x$, choose integer $z_h$ with $hz_h\to x$ as
$h\downarrow0$. Continuity gives $N_k(x,0)=0$. Thus $dA_k=0$.
The radial identity and Euler's identity give, componentwise,
\[
 \sum_i x_i(\partial_i A_{k,j}-\partial_j A_{k,i})
                                      =(k+2)A_{k,j}.
\]
Indeed the first term is $(k+1)A_{k,j}$, and differentiation of
$\sum_i x_iA_{k,i}=0$ makes the second term $-A_{k,j}$.
Hence $A_k=0$. Descending induction proves $A=0$.

A finite subcollection already separates. Start with the whole
space as the common kernel. If that kernel is nonzero, one of the
samplers is nonzero on a vector in it; adding that sampler strictly
reduces the kernel dimension. After finitely many steps the kernel
is zero. Fix the resulting slots $s_1,\ldots,s_m$ and the analytic
map $S=(S_{s_1},\ldots,S_{s_m})$ into the panel of $m$ Lie-algebra
values. Its derivative $D=DS(0)$ is injective.

\subsection{Equivariant recovery and polynomial substitution}
The panel has its invariant Euclidean inner product. Constant gauge
transformations act orthogonally there, preserve $\operatorname{ran}D$,
and intertwine $D$ with their action on $\mathcal J_d$.
Let $\pi$ be orthogonal projection onto $\operatorname{ran}D$ and
put $C=D^{-1}\pi$. The range and its orthogonal complement are
invariant, so $C$ is equivariant and $CD=I$.
The analytic inverse function theorem applied to $CS$ at zero gives
an analytic local inverse $g$. Set $F=gC$. Then
\[
                     F(0)=0,\qquad F(S(A))=A
\]
near zero. For each fixed gauge transformation, equivariance of
$CS$ and uniqueness of its local inverse imply equivariance of
$F$ as germs. This assertion needs no common invariant neighborhood
for all transformations simultaneously.

Let $T_dF$ be the Taylor polynomial of this one recovery map through
ordinary degree $d$. It is equivariant: differentiate the germ identity
$F(Rx)=R'F(x)$ at zero, and use linearity of $R,R'$ in each Taylor
coefficient. It has zero constant term. For each real homogeneous
basis label $p_\alpha$ define a polynomial in panel coordinates by
\[
                            P_\alpha(x)=p_\alpha(T_dF(x)).
\]
It is gauge invariant and real valued on real panels, and $P_0=1$.
If $p_\alpha$ has $r_\alpha$ curvature factors, every monomial of
$P_\alpha$ has ordinary panel degree at least $r_\alpha$:
each substituted coordinate has zero constant term, so each of the
$r_\alpha$ factors contributes at least one panel coordinate.
No deletion of higher-degree terms after substitution is necessary.

These polynomials read the correct jets. Along the radial curve
$A(h)=h\cdot J$, put $W(h)=S(A(h))$. It is smooth and $W(0)=0$.
The Taylor polynomial $T_dF$ and $F$ have identical derivatives
through order $d$ at zero, so their compositions with $W$ do too.
Since $F(W(h))=A(h)$ near zero, composition with the polynomial
$p_\alpha$ gives
\[
 P_\alpha\bigl(s_1(H(J,h)),\ldots,s_m(H(J,h))\bigr)
           \quad\hbox{and}\quad p_\alpha(h\cdot J)
\]
the same derivatives through order $d$.
Linearly extending the $P_\alpha$ therefore gives a local section
$T_0$ of every label in $V_d$ at the level of test jets.

\subsection{Lattice symmetry and a common stencil}
Let $\mathcal H$ be the finite group of signed coordinate
permutations. For its action $\rho$ on labels, define
\begin{equation}\label{eq:local-section-average}
     T(v)(U)=\frac1{|\mathcal H|}\sum_{s\in\mathcal H}
                       T_0(\rho(s^{-1})v)(\operatorname{pull}_sU).
\end{equation}
The test holonomy transforms by the inverse tensor action.
This can be checked directly in the link equation. Positive
coordinate changes transport its coefficient. For a reversed link,
if $u'=uA$, then $u(1)^{-1}u(1-h)$ solves the reversed equation
with coefficient $-A(1-h)$ and starts at the identity. Uniqueness
identifies its endpoint with the inverse link. Constant conjugation
likewise conjugates the solution of the equation. These identities
prove the stated tensor and gauge transformation laws of the sampler.

It follows that each summand of \eqref{eq:local-section-average}
has label $v$. Their average does also. Reindexing the finite group
proves the signed-permutation covariance
$T(v)(\operatorname{pull}_sU)=T(\rho(s)v)(U)$. Its unit
is one, and it commutes with ordinary conjugation because the basis
labels and the initial panel polynomials are real. Each transported
plaquette transforms into another transported plaquette, possibly
with a minus sign from reversing its orientation. Take the finite
union of the slots in all these terms and enclose their paths and
plaquettes in a centered box $Y$.

The average remains polynomial in this common finite list of slots.
Replacing a slot by a signed slot changes no monomial degree.
Moreover $\rho(s)$ preserves the number of curvature factors in
a homogeneous label. Thus the polynomial representing the column
$T(e_\alpha)$ still has no terms of panel degree below $r_\alpha$.
Each slot is itself a polynomial in the positive-link matrix entries
and their conjugates: use matrix multiplication for the path and
plaquette, conjugate transpose for inverse links, and the real-linear
Lie projection of $(M-M^*)/2$ for curvature. A real-linear matrix
functional is a linear combination of real and imaginary entry
parts, each a linear combination of entries and their conjugates.
Therefore $T$ is a genuine finite polynomial cylinder on $Y$.

\subsection{Threshold prices at every period}
Fix derivative order $q$ and spatial power $p$. For any period,
anchor $o$ and allowed direction $A$, a link from the fixed stencil
satisfies $|A(e)|\le C_Yt$, since its torus distance from $o$ is
at most the size of its chosen covering-space position. An ordered
left derivative of a positive-link matrix for the list
$(A_1,\ldots,A_k)$ is exactly
$\operatorname{col}(A_k(e))\cdots\operatorname{col}(A_1(e))U_e$,
where $\operatorname{col}$ is the Lie-algebra matrix representation. The inverse link is its conjugate transpose.
Consequently all derivatives of order $k\le q$ of any fixed path
matrix are bounded by $C t^k$, uniformly in the period and $t$.
The product rule has finitely many terms, and compactness of the
matrix group bounds the undifferentiated factors. Smoothness holds
globally in exponential coordinates.

The same derivative bound holds for each transported plaquette
slot. At order zero a stronger estimate comes from health:
on a configuration healthy on the projected stencil, the plaquette
curvature has norm at most $t$, and parallel transport preserves
that norm. At every positive order $k$, the previous bound is at
most $Ct$ because $0<t\le1$. Thus each slot and each of its tested
derivatives through order $q$ has norm at most $C t$ on that
healthy set. This also holds when the stencil wraps: the path
identities and the bound on torus distance still apply.

Apply the ordered product rule to a monomial with $r$ slot factors.
Each term assigns the derivatives to those factors in their original
orders. Every factor, including those receiving no derivative,
contributes a bound $C_it$. The monomial is therefore bounded by
$C t^r$ through order $q$. Finite summation and $t\le1$ give
\begin{equation}\label{eq:local-section-column-price}
            \|T(e_\alpha)\|_{o,Y}\le C_{d,q,p}t^{r_\alpha}.
\end{equation}
This proves both regularity and the exact factor-count threshold
power of the constructed section. Multiplying its column
coefficients by $t^{-r_\alpha}$ and using the triangle inequality
yields \eqref{eq:local-section-price} with $w=C_{d,q,p}$.

All choices are independent of the period and threshold. They can
also be priced before choosing a scale: for an integer $D$, choose
one bound for each triple $d,q,p\le D$ and set
\[
     w_D=1+\sum_{d,q,p\le D}\max(0,C_{d,q,p}).
\]
It bounds every requested price in this finite range. The stencil
fits whenever the side exceeds
$\max_{x\in Y,i}20(|x_i|+1)$. These are fixed local requirements
for the later calendar to satisfy. The construction here supplies
the section and its uniform price; it does not choose that calendar.
